\section*{Chapter \ref{ch:m2:bond-futures} --- Bond Futures}

\begin{solution}{exo:m2:bond-futures:1}
$111.765625 \times 0.8870 + 1.453125 = 100.589234$ per 100; USD~100\,589.23 per
contract of USD~100\,000 face.
\end{solution}

\begin{solution}{exo:m2:bond-futures:2}
$99.132 - 111.765625 \times 0.8870 = -0.004$, about $-0.14$ of a 32nd. The \omterm{def:m2:bond-futures:irr}{gross
basis} compares today's spot price with the futures price scaled by the factor;
the future is a forward, so the gap also contains the carry to delivery.
Here carry is negative (the note's coupon, 3.875\%, earns less than 3.90\% repo
on its price), so the forward is above spot and the \omterm{def:m2:bond-futures:irr}{gross basis} can be
negative while the \omterm{def:m2:bond-futures:irr}{net basis} is positive.
\end{solution}

\begin{solution}{exo:m2:bond-futures:3}
From 1 December 2026 to 15 November 2033: 6 years 11 months, rounded down to
6 years 9 months, so $n = 6$, $z = 9$, $v = 3$. The 4\% note: 0.8902. The 6\%
note: 0.9999, not 1, because the formula prices it as if its next coupon were
three months away and its life rounded down; only a 6\% note whose life is a
whole number of quarters from the first day of the month has a factor of
exactly 1.
\end{solution}

\begin{solution}{exo:m2:bond-futures:4}
Earning 3.8238\% and paying 3.90\% on the dirty price, USD~99\,563\,440, for 97
days loses $99\,563\,440 \times 0.000762 \times 97/360 = \text{USD}~20\,451$, if the
note is delivered as the CTD. One still buys the basis to own the short's options:
if yields move far enough for the CTD to switch, or if the timing and \omterm{def:m2:bond-futures:options}{wildcard
options} pay, the basis widens and the long profits; and if the note becomes
special in repo, its financing costs less than 3.90\%.
\end{solution}

\begin{solution}{exo:m2:bond-futures:5}
The future moves like the CTD divided by its factor, so its DV01 per contract
is the CTD's DV01 per USD~100\,000 divided by the factor. Contracts
$= \frac{100\,000\,000/100\,000 \times \mathrm{DV01}_{\text{CTD}}}{\mathrm{DV01}_{\text{CTD}}/\mathrm{CF}} = 1\,000 \times 0.8870 = 887$
contracts. The factor is the number of futures ``per bond''.
\end{solution}

\begin{solution}{exo:m2:bond-futures:6}
The factor prices every note at 6\%. Below 6\% all notes are worth more than
their factor-scaled value, and the longer the duration, the more they have
gained as yields fell from 6\%: the long notes are the dearest to deliver, the
shortest the cheapest. As yields rise towards 6\% the long notes lose more,
the gap shrinks, and at some yield (below 6\% here, because the curve and the
coupons differ) the longest note becomes the cheaper one.
\end{solution}

\begin{solution}{exo:m2:bond-futures:7}
The standard deviation over 97 days is $80 \times \sqrt{97/365} = 41$ basis points.
The switch option is worth 0.001 of a 64th: the switch is 3.6 standard
deviations away. At three times the volatility, 124 basis points, it is worth
3.07 64ths. The \omterm{def:m2:bond-futures:options}{quality option} is worth nothing until the switch comes within
reach, then grows fast: it is an out-of-the-money option on the level of
yields.
\end{solution}

\begin{solution}{exo:m2:bond-futures:8}
The future is priced per 100 of a notional 6\% note; the deliverable notes pay
about 4\% and are worth about 100 at 4\% yields. A 6\% coupon note of the same
maturity would be worth about 112 at those yields. The factors (0.88 to 0.92)
convert: $111.77 \times 0.887 = 99.14$, within a hundredth of the CTD's 99.13. The comparison must be
made after multiplying by the factor, where the gap is a few hundredths of a
point, not 11\%.
\end{solution}

\begin{solution}{pb:m2:bond-futures:1}
\textbf{1.} From 1 December 2026: Aug-33 6 years 8 months, Nov-33 6 years 11
months, Feb-34 7 years 2 months, May-34 7 years 5 months, Aug-34 7 years 8
months: all at least 6 years 6 months and less than 8 years, and all issued as
notes of at most ten years: all deliverable.
\textbf{2.} Aug-34: 0.8806. CTD (Aug-33): 0.8870.
\textbf{3.} \omterm{def:m2:bond-futures:irr}{Gross basis} $-0.004$; carry $-0.025$; \omterm{def:m2:bond-futures:irr}{net basis} $-0.004 + 0.025 = 0.020$.
\textbf{4.} 3.82\%, eight basis points below repo.
\textbf{5.} Its coupon income over the 97 days, 1.021 per 100, is less than the
repo interest on its dirty price, 1.046: a coupon of 3.875\% a year on face
accrues more slowly than 3.90\% charged on a 360-day year.
\textbf{6.} Nov-33 0.270, Feb-34 0.527, May-34 0.742, Aug-34 1.171 points above the
CTD's 111.793.
\textbf{7.} $+147$ basis points; the 4\% note of August 2034, the longest,
takes over.
\textbf{8.} $4.02\% + 1.47\% = 5.49\%$ (the Aug-34 then yields 5.55\%).
\textbf{9.} The notes do not all yield the same: the longer ones yield more, as
the curve slopes up, and they differ in coupon. The switch happens where the
longer note's extra yield compensates the factor's bias, below the 6\% where
a flat curve would put it.
\textbf{10.} The future's DV01 follows the new CTD's DV01 per factor: from about
586 to 658 per 100 of futures notional (in units of $10^{-4}$), 12\% more. A
hedger who holds a fixed number of contracts is suddenly over-hedged.
\textbf{11.} $80\sqrt{97/365} = 41$ basis points.
\textbf{12.} $147/41 = 3.6$ standard deviations.
\textbf{13.} 0.001 of a 64th: worthless.
\textbf{14.} 3.07 64ths.
\textbf{15.} The timing and \omterm{def:m2:bond-futures:options}{wildcard options}, the end-of-month option, and the
market's uncertainty about the \omterm{def:m2:repo-and-specials:rate}{repo rate} the CTD will finance at: the CTD may
turn special and its forward be lower than the general-collateral forward.
\textbf{16.} Long the basis she is long the note and short the future, which
is short the options the futures short owns: she owns them. Options gain from
volatility: a large move either way can make another note cheaper, and her
long note's price then rises relative to the future.
\textbf{17.} Because the future's duration jumps when the CTD changes: a hedge
ratio computed from the old CTD is wrong after the switch, just when yields
have moved a lot.
\textbf{18.} More supply of the CTD: it cheapens in the cash market and in repo
(less special), its \omterm{def:m2:bond-futures:irr}{net basis} falls, and it becomes even more firmly the
cheapest to deliver.
\textbf{19.} Named result: \emph{the switch point} is a parallel rise of 147
basis points (CTD yield 5.49\%), when the 4\% of August 2034 replaces the
3.875\% of August 2033; the switch option is worth 0.001 of a 64th at 80 basis
points of annual volatility and 3.07 64ths at 240.
\textbf{20.} The cheapest note in its basket, divided by its factor, less the
value of the short's choices.
\end{solution}

\begin{solution}{iq:m2:bond-futures:1}
The contract allows several notes with different coupons and maturities to be
delivered. The factor, the note's price at a 6\% yield, scales the invoice so
that each note is delivered at a price proportional to its value if yields
were 6\%. Without it the short would always deliver the lowest-priced note and
the contract would be written on one issue, open to squeezes.
\iqlookfor{what the factor is (a 6\% price) and why a basket needs it.}
\end{solution}

\begin{solution}{iq:m2:bond-futures:2}
For each deliverable, compute the \omterm{def:m2:bond-futures:irr}{implied repo rate} from its price, the futures
price, the factor, accrued interest and coupons to delivery; the highest is
the CTD. Equivalently, the lowest \omterm{def:m2:bond-futures:irr}{net basis}. At delivery, the lowest price
divided by factor.
\iqlookfor{implied repo, not the lowest price or lowest \omterm{def:m2:bond-futures:irr}{gross basis}.}
\end{solution}

\begin{solution}{iq:m2:bond-futures:3}
The financing rate at which buying the note, financing it and delivering it
into the future breaks even. If the CTD's implied repo exceeds actual repo,
buying the basis and delivering locks in a profit: the future is rich. In
practice the CTD's implied repo sits a little below repo, the difference paying
for the short's delivery options.
\iqlookfor{the arbitrage interpretation and why the normal state is below
repo.}
\end{solution}

\begin{solution}{iq:m2:bond-futures:4}
Buy the CTD, finance it in repo, sell futures in the ratio of the factor;
earn the \omterm{def:m2:bond-futures:irr}{implied repo rate} against the \omterm{def:m2:repo-and-specials:rate}{repo rate} paid, a few basis points,
levered many times through repo. What can go wrong: \omterm{def:m2:repo-and-specials:rate}{repo rates} spike or
lenders raise haircuts, so financing costs more or must be reduced; futures
margin calls on the short side during a rally must be met in cash while the
gain on the notes is unrealised; the basis widens in a \omterm{def:m2:the-treasury-market:dash}{dash for cash}; and many
leveraged funds unwinding at once can move the market they unwind in.
\iqlookfor{the financing and margin mechanics, and the funding liquidity
risk, not only the arbitrage.}
\end{solution}

\begin{solution}{iq:m2:bond-futures:5}
To first order, the CTD's DV01 per 100 of face divided by its \omterm{def:m2:bond-futures:cf}{conversion
factor}, times 1\,000 for a USD~100\,000 contract. Better: reprice the future
from the whole basket under bumped yields (the minimum of forward prices per
factor, plus an option model), which captures the change in CTD and the
optionality; the two agree only far from a switch.
\iqlookfor{DV01 over factor, and why the switch makes it state-dependent.}
\end{solution}

\begin{solution}{iq:m2:bond-futures:6}
The \omterm{def:m2:bond-futures:options}{quality option} (which note), the timing option (which day of the delivery
month), the \omterm{def:m2:bond-futures:options}{wildcard option} (deliver after the futures price is fixed for the
day, or after trading ends at a fixed \omterm{def:m2:bond-futures:cf}{invoice price}), and the end-of-month
option. The \omterm{def:m2:bond-futures:options}{quality option} is usually the largest: value it by simulating the
yield curve to delivery (a factor model of parallel, slope and curvature moves),
computing each note's price per factor, and taking the expected minimum
against the current CTD's; volatility and the distance to the switch drive it.
\iqlookfor{naming the options and a concrete method for the \omterm{def:m2:bond-futures:options}{quality option}.}
\end{solution}
